%--- E8: Section 5.3, after Eq. (lagrange) ----------------------------------
\begin{explained}{the Lagrange projector for evolution factors}
\textbf{Step 1: the evolution operator.} $P(\tau,\xi)=\prod_{j=1}^N(\tau-a_j(\xi))$ with $a_j$ distinct polynomials in $\xi$. Replacing $\tau,\xi$ by $\partial_t,\partial_x$ the equation is $\prod_j(\partial_t-a_j(\partial_x))u=0$. Example: $a_1(\xi)=\xi$, $a_2(\xi)=-\xi$ gives $(\partial_t-\partial_x)(\partial_t+\partial_x)u=0$, the one-dimensional wave equation $u_{tt}-u_{xx}=0$. Because the factors are distinct, a solution decomposes as $u=\sum_ju_j$ where each component satisfies its own first-order law $\partial_tu_j=a_j(\partial_x)u_j$, i.e.\ $(\partial_t-a_j(\partial_x))u_j=0$.

\medskip\textbf{Step 2: isolating a component.} Fix $j$ and apply $Q_j=\prod_{k\ne j}(\partial_t-a_k(\partial_x))$ to $u=\sum_lu_l$. Constant-coefficient operators commute, so $Q_ju=\sum_lQ_ju_l$.
\begin{itemize}[nosep]
\item \emph{Case $l\ne j$.} $Q_j$ contains the factor $\partial_t-a_l(\partial_x)$, which kills $u_l$: $Q_ju_l=0$. Every component except $u_j$ disappears.
\item \emph{Case $l=j$.} No factor of $Q_j$ kills $u_j$; use $\partial_tu_j=a_j(\partial_x)u_j$: for each $k\ne j$, $(\partial_t-a_k(\partial_x))u_j=(a_j-a_k)(\partial_x)u_j$. Applying all factors, $Q_ju=\prod_{k\ne j}(a_j-a_k)(\partial_x)\,u_j$.
\end{itemize}
Hence
\[
u_j=\Big[\prod_{k\ne j}(a_j-a_k)(\partial_x)\Big]^{-1}\prod_{k\ne j}(\partial_t-a_k(\partial_x))\,u,
\]
which is Eq.~\eqref{eq:lagrange}, provided the inverse operator exists on the function space considered. It is not automatically a differential operator: it is in general a Fourier multiplier, possibly singular at frequencies where $a_j=a_k$.

\medskip\textbf{Step 3: why this is Lagrange interpolation.} The Lagrange polynomials $L_j(z)=\prod_{k\ne j}\frac{z-a_k}{a_j-a_k}$ satisfy $L_j(a_l)=1$ for $l=j$, $0$ otherwise; for a matrix $C$ diagonalized in eigenvectors, $L_j(C)$ acts as the identity on the $j$-th eigenspace and as zero on the others. Here the role of $C$ is played by $\partial_t$ and the eigenvalues $a_j$ are replaced by the operators $a_j(\partial_x)$:
\[
\Pi_j=\prod_{k\ne j}\frac{\partial_t-a_k(\partial_x)}{a_j(\partial_x)-a_k(\partial_x)},\qquad \Pi_ju=u_j .
\]
The analogy is exact at the algebraic level: $\Pi_j$ is the projector that extracts the $j$-th evolution component.

\medskip\textbf{Step 4: rational in $\xi$.} With the Fourier convention $\partial_x\mapsto i\kappa$ ($\kappa$ real frequency), $a_j(\partial_x)$ becomes multiplication by $a_j(i\kappa)$ and the multiplier is
\[
\ell_j(\kappa,\partial_t)=\prod_{k\ne j}\frac{\partial_t-a_k(i\kappa)}{a_j(i\kappa)-a_k(i\kappa)} .
\]
The denominators vanish exactly where $a_j(i\kappa)=a_k(i\kappa)$: these are the potential poles of the projector. Strictly, the rationality concerns the dependence on the spatial frequency; the projector still involves $\partial_t$ before it is restricted to a solution component, on which $\partial_t$ acts as multiplication by $a_j(i\kappa)$.

\medskip\textbf{Step 5: the wave equation.} $(\partial_t-\partial_x)(\partial_t+\partial_x)u=0$, $u=u_1+u_2$ with $\partial_tu_1=\partial_xu_1$ and $\partial_tu_2=-\partial_xu_2$. Applying $\partial_t+\partial_x$: the $u_2$ term vanishes and $(\partial_t+\partial_x)u_1=2\partial_xu_1$, so
\[
u_1=\tfrac12\,\partial_x^{-1}(\partial_t+\partial_x)u,\qquad u_2=-\tfrac12\,\partial_x^{-1}(\partial_t-\partial_x)u .
\]
In Fourier variables $\partial_x^{-1}$ is multiplication by $1/(i\kappa)$: a pole at $\kappa=0$. This does not mean the wave equation has no decomposition (Section~\ref{sec:fact} proves it does), only that this particular formula needs a choice of inverse (an integration constant) or extra assumptions at zero frequency.

\medskip\textbf{Step 6: what ``coprime'' means.} $\tau-a_j(\xi)$ and $\tau-a_k(\xi)$ are coprime in $\C[\tau,\xi]$ exactly when $a_j-a_k$ is not the zero polynomial; for distinct $a_j$ this holds, though $a_j-a_k$ may still vanish at particular frequencies, creating poles. Three levels must be kept apart: \emph{coprime factors} (the algebraic decomposition exists in the localized, rational-operator sense); \emph{no poles at any frequency} (stronger, not guaranteed by coprimality); \emph{a bounded or well-defined projector on a chosen function space} (needs analytic conditions near frequencies where branches meet). The message: distinct evolution branches admit algebraic projectors, like distinct eigenvalues of a matrix; in spatial Fourier variables the projectors are rational multipliers whose denominators reveal where the branches intersect.
\end{explained}

